%%%%%%%%%%%%%%%%%%%%%%%%%%%%%%%%%%%%%%%%%%%% SIMBOLOS %%%%%%%%%%%%%%%%%%%%%%%%%%%%%%%%%%%%%%%%%%%%%%%%%%%%

\renewcommand{\qedsymbol}{\textcolor{rojoCIMAT}{\ensuremath{\square}}}
% Redefinir \mathbb para usar Latin Modern
\DeclareMathAlphabet{\mathbb}{U}{lmsb}{m}{n}
% Redefinimos partes real e imaginaria
\renewcommand{\Re}{\mathcal{R}\textit{e}}
\renewcommand{\Im}{\mathcal{I}\textit{m}}
% Agregamos alternativas reescalables de \ulcorner y \urcorner del paquete MnSymbols
\makeatletter
\DeclareFontFamily{OMX}{MnSymbolE}{}
\DeclareSymbolFont{MnLargeSymbols}{OMX}{MnSymbolE}{m}{n}
\SetSymbolFont{MnLargeSymbols}{bold}{OMX}{MnSymbolE}{b}{n}
\DeclareFontShape{OMX}{MnSymbolE}{m}{n}{
    <-6>  MnSymbolE5
   <6-7>  MnSymbolE6
   <7-8>  MnSymbolE7
   <8-9>  MnSymbolE8
   <9-10> MnSymbolE9
  <10-12> MnSymbolE10
  <12->   MnSymbolE12
}{}
\DeclareFontShape{OMX}{MnSymbolE}{b}{n}{
    <-6>  MnSymbolE-Bold5
   <6-7>  MnSymbolE-Bold6
   <7-8>  MnSymbolE-Bold7
   <8-9>  MnSymbolE-Bold8
   <9-10> MnSymbolE-Bold9
  <10-12> MnSymbolE-Bold10
  <12->   MnSymbolE-Bold12
}{}
\DeclareMathDelimiter{\ulcor}
    {\mathopen}{MnLargeSymbols}{'036}{MnLargeSymbols}{'036}
\DeclareMathDelimiter{\urcor}
    {\mathclose}{MnLargeSymbols}{'043}{MnLargeSymbols}{'043}
\makeatother

\newcommand{\bigforall}{\scalebox{1.2}{$\forall$}} % Simbolo para todo grande
\newcommand{\bigexists}{\scalebox{1.2}{$\exists$}} % Simbolo existe grande
% Símbolo pertenencia
\newcommand{\en}{\mathrel{% 
    \mathchoice%
    {\hspace{-0.2em}\raisebox{0.1ex}{\scalebox{0.95}{$\in$}}\hspace{-0.2em}} % Tamaño Displaymath
    {\hspace{-0.2em}\raisebox{0.1ex}{\scalebox{0.95}{$\in$}}\hspace{-0.2em}} % Tamaño Inline
    {\hspace{-0.03em}\raisebox{0.07ex}{\scalebox{0.65}{$\in$}}\hspace{-0.03em}} % Scriptstyle
    {\hspace{-0.03em}\raisebox{0.05ex}{\scalebox{0.6}{$\in$}}\hspace{-0.03em}} % Scriptscriptstyle
}}
\newcommand{\noten}{\mathrel{% 
    \mathchoice%
    {\hspace{-0.2em}\raisebox{0.1ex}{\scalebox{0.95}{$\notin$}}\hspace{-0.2em}} % Tamaño Displaymath
    {\hspace{-0.2em}\raisebox{0.1ex}{\scalebox{0.95}{$\notin$}}\hspace{-0.2em}} % Tamaño Inline
    {\hspace{-0.03em}\raisebox{0.07ex}{\scalebox{0.65}{$\notin$}}\hspace{-0.03em}} % Scriptstyle
    {\hspace{-0.03em}\raisebox{0.05ex}{\scalebox{0.6}{$\notin$}}\hspace{-0.03em}} % Scriptscriptstyle
}}
\newcommand{\cont}{\mathrel{% 
    \mathchoice%
    {\hspace{-0.2em}\raisebox{0.1ex}{\scalebox{0.95}{$\ni$}}\hspace{-0.2em}} % Tamaño Displaymath
    {\hspace{-0.2em}\raisebox{0.1ex}{\scalebox{0.95}{$\ni$}}\hspace{-0.2em}} % Tamaño Inline
    {\hspace{-0.03em}\raisebox{0.07ex}{\scalebox{0.65}{$\ni$}}\hspace{-0.03em}} % Scriptstyle
    {\hspace{-0.03em}\raisebox{0.05ex}{\scalebox{0.6}{$\ni$}}\hspace{-0.03em}} % Scriptscriptstyle
}}
\newcommand{\supsub}[3]{{#1}^{\scriptscriptstyle #2}_{\scriptscriptstyle #3}} % Comando exponentes más chicos
\newcommand{\restr}[2]{{#1}{\left|_{ {#2}}\kern-\nulldelimiterspace\right.}} % Restricción extensible
% Comando barra superior extensible 
\newcommand{\ubar}[1]{\mathpalette\ubaraux{#1}}

\newcommand{\ubaraux}[2]{%
  \tikz[baseline=(X.base)]{
    \node[inner sep=0pt, outer sep=0pt] (X) {$#1#2$};
    \draw[line width=0.5pt]
      ([xshift=1pt, yshift=1pt]X.north west) --
      ([xshift=-1pt, yshift=1pt]X.north east);
  }%
}
% Definición delimitadores
\providecommand\given{}% Aseguramos que \given exista
\newcommand\SetSymbol[1][]{% Símbolo "|" con soporte par múltiple línea
    \nonscript\:#1\vert
    \allowbreak
    \nonscript\:
    \mathopen{}
}
\DeclarePairedDelimiterX\ang[1]{\langle}{\rangle}{\renewcommand\given{\SetSymbol[\delimsize]}#1}% <, > reescalables 
\DeclarePairedDelimiterX\set[1]\{\}{\renewcommand\given{\SetSymbol[\delimsize]}#1}% {, } para conjuntos
\DeclarePairedDelimiter\cornered\ulcor\urcor% Esquinas superiores como delimitadores reescalables
\DeclarePairedDelimiter\Brack\lBrack\rBrack
\DeclarePairedDelimiter\round\lfloor\rceil
\let\brack\relax
\DeclarePairedDelimiter\brack{[}{]}
\DeclarePairedDelimiter\floor\lfloor\rfloor
\DeclarePairedDelimiter\ceil\lceil\rceil

% Definición operadores útiles
\DeclareMathOperator{\Arg}{Arg}
\DeclareMathOperator{\supp}{supp}
\DeclareMathOperator{\diam}{diam}
\newcommand{\dotcup}{\mathbin{\dot{\cup}}}
\newcommand{\ctpeq}{\xlongequal{\hspace{-0.3em}\text{a.e.}\hspace{-0.3em}}}
\DeclareMathOperator{\vol}{Vol}
\DeclareMathOperator{\gal}{Gal}
\DeclareMathOperator{\p}{\textit{p}}
\DeclareMathOperator{\rk}{rk}
\DeclareMathOperator{\card}{card}
\DeclareMathOperator{\inter}{int}
\DeclareMathOperator{\aut}{Aut}
\DeclareMathOperator{\fv}{FV}
\DeclareMathOperator{\coker}{coker}
\DeclareMathOperator{\im}{im}
\DeclareMathOperator{\pert}{\ni}
\DeclareMathOperator{\desc}{Desc}
\DeclareMathOperator{\dom}{dom}
\DeclareMathOperator{\cod}{cod}
\DeclareMathOperator{\sv}{Sv}
\DeclareMathOperator{\et}{\mathscr{E}\!}
\DeclareMathOperator{\sh}{Sh}
\DeclareMathOperator{\spsh}{SepPSh}
\DeclareMathOperator{\psh}{PSh}
\DeclareMathOperator{\res}{res}
\DeclareMathOperator{\Spec}{Spec}
\DeclareMathOperator{\Open}{Open}
\DeclareMathOperator{\true}{\top}
\DeclareMathOperator{\false}{\bot}
\let\ev\relax
\DeclareMathOperator{\ev}{\textit{ev}}
\DeclareMathOperator{\ob}{Ob}
\DeclareMathOperator{\id}{id}
\DeclareMathOperator{\sub}{Sub}
\DeclareMathOperator{\cone}{Cone}
\DeclareMathOperator{\cocone}{Cocone}
\DeclareMathOperator*{\colim}{colim} % Operador colímite sin acento y se comporta bien con subindices
\newcommand{\Col}[1]{\operatorname{#1}} % Comando genérico para colecciones de objetos
\newcommand{\cat}[1]{\textit{#1}}
% Categorías con el estilo del proyecto de apuntes.
\newcommand{\Ring}{\mathbf{Ring}}
\newcommand{\CRing}{\mathbf{CRing}}
\newcommand{\AffSch}{\mathbf{AffSch}}
\newcommand{\Sch}{\mathbf{Sch}}
\newcommand{\LocAn}{\mathbf{LocAn}}
\newcommand{\coamalg}{\rotatebox[origin=c]{180}{$\amalg$}} % Simbolo producto binario alternativo
\renewcommand{\hom}{\operatorname{Mor}} % Redefinimos \hom 
\DeclareMathOperator{\Hom}{Hom}\newcommand{\cech}[1]{\operatorname{\check{H}^{#1}}}
\newcommand{\colonequiv}{\vcentcolon\equiv}

%%%%%%%%%%%%%%%%%%%%%%%%%%%%%%%%%%%%%%%%%%%% ENTORNOS %%%%%%%%%%%%%%%%%%%%%%%%%%%%%%%%%%%%%%%%%%%%%%%%%%%%

\declaretheorem[name = Teorema]{teo}
\crefname{teo}{teorema}{teoremas}
% Subentornos de teo que se nombren adecuadamente con \cref
\declaretheorem[name = Proposición, sibling = teo]{prop}
\crefname{prop}{proposición}{proposiciones}

\declaretheorem[name = Lema, sibling = teo]{lema}
\crefname{lema}{lema}{lemas}

\declaretheorem[name = Corolario, sibling = teo]{coro}
\crefname{coro}{corolario}{corolarios}

% Otros entornos matemáticos
% \theoremstyle{definition}
% \declaretheorem[name = Definición, numberwithin = section]{defi}
% \declaretheorem[name = Ejemplo, numberwithin = section]{ejem}

% Entorno auxiliar para ejercicio
\declaretheorem[name = Ejercicio, numberwithin = section]{eje}
\declaretheorem[name = Problema, numberwithin = subsection]{pro}
\declaretheorem[name = Ejercicio, numberwithin = subsection]{neje}
\crefname{eje}{ejercicio}{ejercicios}
\crefname{neje}{ejercicio}{ejercicios}
\crefname{pro}{problema}{problemas}

% Cambiamos formato de contadores
\renewcommand{\theneje}{\arabic{section}.\arabic{subsection}.\Alph{neje}}
\renewcommand{\theHneje}{%
    \arabic{section}.\arabic{subsection}.\arabic{neje}%
} % Cambiamos formato del ancla de hyperref, de otro modo superpone etiquetas entre subsecciones
\renewcommand{\theeje}{\arabic{eje}}

\makeatletter
% Entorno ejercicio que se agrega al TOC 
\newenvironment{ejer}[1][]{%
    \refstepcounter{eje} %
    \label{ejer:\thechapter.\theeje} % Etiqueta automática para referencias
    \if\relax\detokenize{#1}\relax
        \def\ejetitle{Ejercicio \theeje}%
        \addcontentsline{toc}{subsubsection}{\ejetitle}%
    \else
        \def\ejetitle{Ejercicio \theeje\ (#1)}%
        \addcontentsline{toc}{subsubsection}{\ejetitle}%
    \fi
    \begin{tcolorbox}[
        enhanced,
        breakable,
        before skip=12pt,
        after skip=12pt,
        top = 0.1em,
        coltitle=rojoCIMAT,
        colback=red!0!rojoCIMAT!5!white,
        colframe=rojoCIMAT,
        colbacktitle=rojoCIMAT,
        boxrule=0.8pt,
        title=\itshape\ejetitle,
        attach boxed title to top center={yshift*=1mm-\tcboxedtitleheight},
        % attach boxed title to top left={xshift=1cm,yshift*=1mm-\tcboxedtitleheight}, % Alternativa pestana a la izquierda
        varwidth boxed title*=-3cm,
        boxed title style={%
            top=-0.05em, % Espacio sobre titulo
            bottom=-0.15em, % Espacio debajo del titulo 
            frame code={%
                \path[draw=tcbcolback,line width=0.8pt]
                ([yshift=-1mm,xshift=-1.1mm]frame.north west)
                arc[start angle=0,end angle=180,radius=1mm]
                ([yshift=-1mm,xshift=1.1mm]frame.north east)
                arc[start angle=180,end angle=0,radius=1mm];
                \path[left color=tcbcolback!2.5,right color=tcbcolback!2.5,middle color=tcbcolback!2.5]
                ([xshift=-2mm]frame.north west)--([xshift=2mm]frame.north east)
                [rounded corners=1mm]--([xshift=1mm,yshift=-1mm]frame.north east)
                --(frame.south east)--(frame.south west)
                --([xshift=-1mm,yshift=-1mm]frame.north west)
                [sharp corners]--cycle;
                \path[draw=tcbcolback,line width=0.8pt]
                ([xshift=-2mm]frame.north west)--([xshift=2mm]frame.north east)
                [rounded corners=1mm]--([xshift=1mm,yshift=-1mm]frame.north east)
                --(frame.south east)--(frame.south west)
                --([xshift=-1mm,yshift=-1mm]frame.north west)
                [sharp corners]--cycle;
            },
            interior engine=empty%
        },
        fonttitle=\bfseries%
    ]%
}{%
    \par\nointerlineskip\vspace{0.1em}%
    \end{tcolorbox}%
}
% Entorno problema que se agrega al TOC 
\newenvironment{prob}[1][]{%
    \refstepcounter{pro} %
    \label{prob:\thepro} % Etiqueta automática para referencias
    \if\relax\detokenize{#1}\relax
        \def\protitle{Problema \thepro}%
        \addcontentsline{toc}{subsubsection}{\protitle}%
    \else
        \def\protitle{Problema \thepro\ (#1)}%
        \addcontentsline{toc}{subsubsection}{\protitle}%
    \fi
    \begin{tcolorbox}[
        enhanced,
        breakable,
        before skip=12pt,
        after skip=12pt,
        top = 0.1em,
        coltitle=rojoCIMAT,
        colback=red!0!rojoCIMAT!5!white,
        colframe=rojoCIMAT,
        colbacktitle=rojoCIMAT,
        boxrule=0.8pt,
        title=\itshape\protitle,
        attach boxed title to top center={yshift*=1mm-\tcboxedtitleheight},
        % attach boxed title to top left={xshift=1cm,yshift*=1mm-\tcboxedtitleheight}, % Alternativa pestana a la izquierda
        varwidth boxed title*=-3cm,
        boxed title style={%
            top=-0.05em, % Espacio sobre titulo
            bottom=-0.15em, % Espacio debajo del titulo 
            frame code={%
                \path[draw=tcbcolback,line width=0.8pt]
                ([yshift=-1mm,xshift=-1.1mm]frame.north west)
                arc[start angle=0,end angle=180,radius=1mm]
                ([yshift=-1mm,xshift=1.1mm]frame.north east)
                arc[start angle=180,end angle=0,radius=1mm];
                \path[left color=tcbcolback!2.5,right color=tcbcolback!2.5,middle color=tcbcolback!2.5]
                ([xshift=-2mm]frame.north west)--([xshift=2mm]frame.north east)
                [rounded corners=1mm]--([xshift=1mm,yshift=-1mm]frame.north east)
                --(frame.south east)--(frame.south west)
                --([xshift=-1mm,yshift=-1mm]frame.north west)
                [sharp corners]--cycle;
                \path[draw=tcbcolback,line width=0.8pt]
                ([xshift=-2mm]frame.north west)--([xshift=2mm]frame.north east)
                [rounded corners=1mm]--([xshift=1mm,yshift=-1mm]frame.north east)
                --(frame.south east)--(frame.south west)
                --([xshift=-1mm,yshift=-1mm]frame.north west)
                [sharp corners]--cycle;
            },
            interior engine=empty%
        },
        fonttitle=\bfseries%
    ]%
}{%
    \par\nointerlineskip\vspace{0.1em}%
    \end{tcolorbox}%
}
% Entorno ejercicio que se agrega al TOC (con número de sección)
\newenvironment{nejer}[1][]{%
    \refstepcounter{neje} %
    \label{nejer:\theneje} % Etiqueta automática para referencias
    \if\relax\detokenize{#1}\relax
        \def\nejetitle{Ejercicio \theneje}%
        \addcontentsline{toc}{subsubsection}{\nejetitle}%
    \else
        \def\nejetitle{Ejercicio \theneje\ (#1)}%
        \addcontentsline{toc}{subsubsection}{\nejetitle}%
    \fi
    \begin{tcolorbox}[
        enhanced,
        breakable,
        before skip=12pt,
        after skip=12pt,
        top = 0.1em,
        coltitle=rojoCIMAT,
        colback=red!0!rojoCIMAT!5!white,
        colframe=rojoCIMAT,
        colbacktitle=rojoCIMAT,
        boxrule=0.8pt,
        title=\itshape\nejetitle,
        attach boxed title to top center={yshift*=1mm-\tcboxedtitleheight},
        % attach boxed title to top left={xshift=1cm,yshift*=1mm-\tcboxedtitleheight}, % Alternativa pestana a la izquierda
        varwidth boxed title*=-3cm,
        boxed title style={%
            top=-0.05em, % Espacio sobre titulo
            bottom=-0.15em, % Espacio debajo del titulo 
            frame code={%
                \path[draw=tcbcolback,line width=0.8pt]
                ([yshift=-1mm,xshift=-1.1mm]frame.north west)
                arc[start angle=0,end angle=180,radius=1mm]
                ([yshift=-1mm,xshift=1.1mm]frame.north east)
                arc[start angle=180,end angle=0,radius=1mm];
                \path[left color=tcbcolback!2.5,right color=tcbcolback!2.5,middle color=tcbcolback!2.5]
                ([xshift=-2mm]frame.north west)--([xshift=2mm]frame.north east)
                [rounded corners=1mm]--([xshift=1mm,yshift=-1mm]frame.north east)
                --(frame.south east)--(frame.south west)
                --([xshift=-1mm,yshift=-1mm]frame.north west)
                [sharp corners]--cycle;
                \path[draw=tcbcolback,line width=0.8pt]
                ([xshift=-2mm]frame.north west)--([xshift=2mm]frame.north east)
                [rounded corners=1mm]--([xshift=1mm,yshift=-1mm]frame.north east)
                --(frame.south east)--(frame.south west)
                --([xshift=-1mm,yshift=-1mm]frame.north west)
                [sharp corners]--cycle;
            },
            interior engine=empty%
        },
        fonttitle=\bfseries%
    ]
}{%
    \par\nointerlineskip\vspace{0.1em}%  
    \end{tcolorbox}%
}
\makeatother

% Entornos notas
\theoremstyle{remark}
\declaretheorem[name = Observación, numbered = no]{obs}
\declaretheorem[name = Nota, numbered = no]{nota}

% Entorno solucion
\newenvironment{sol}[1][Solución]{%
    \begin{proof}[#1]
}{%
    \end{proof}
    \vspace{1em}
    % \pagebreak[3]
    \clearpage
}

% Referencias inteligentes "\cref (\nameref)"
\newcommand{\myref}[1]{%
    \cref{#1}%
    \begingroup
    \edef\temp{\getrefbykeydefault{#1}{name}{}}%
    \ifx\temp\empty
        % No hacer nada si \temp está vacío
    \else
        ~(\nameref{#1})% Imprime nombre entre paréntesis en caso de existir
    \fi
    \endgroup
}

% Comando para referencias abreviadas
\newcommand{\abvref}[1]{%
    \begingroup
    \crefname{teo}{teo.}{teos.}%
    \crefname{lema}{lema}{lemas}%
    \crefname{prop}{prop.}{props.}%
    \crefname{coro}{coro.}{coros.}%
    \crefname{section}{sec.}{secs.}%
    \crefname{eje}{ej.}{ejs.}%
    \cref{#1}%
    \endgroup
}
